\section{\iflanguage{english}{Fundamentals}{Grundlagen}}\label{sec:fundamentals}

\subsection{The Practical Algebraic Calculus}\label{sec:practical-algebraic-calculus}

This section introduces the practical algebraic calculus (PAC) as it is used by Pastèque.
For a thorough introduction on PAC, I refer to the work of Kaufmann et al. \cite{pac-2018} \cite{pac-2024} \cite{kaufmann-phd-thesis}.

\begin{itemize}
  \item A PAC checker considers multivariate polynomials $\mathbb{Z}[X]$ with variables $X$ and coefficients in $\mathbb{Z}$.

  \item Variables $X$ represent Boolean values; for each $x \in X$ it is required that $x^2 - x = 0$.
        $B(X) = \{x^2 - x \mid x \in X\}$ is called the set of \textit{Boolean value constraints} over the variables $X$.
  \item For a set $P \subseteq \mathbb{Z}[X]$, a \textit{model} is a point $m \in \mathbb{Z}^{\Vert X \Vert}$ such that every $p \in P$ is 0 at $m$, i.e. $\forall p \in P.\, p(m) = 0$.
  \item A nonempty subset $I \subseteq \mathbb{Z}[X]$ is called \textit{ideal} if it is closed under addition and closed under  multiplication by arbitrary polynomials, i.e.
        $I$ is an ideal iff. for all $i_1,i_2\in I$ and $p \in \mathbb{Z}[X]$ it holds that $i_1 + i_2 \in I$ and $p \cdot i_1 \in I$.
  \item For a set $B = \{ b_1, \dots, b_n \} \subseteq \mathbb{Z}[X]$, $\langle B \rangle$ denotes the \textit{ideal generated by} $B$, where $\langle B \rangle = \{b_1q_1 + \dots + b_nq_n \mid q_1,\dots,q_n \in \mathbb{Z}[X]\}$.
\end{itemize}

In the practical algebraic calculus, polynomials model the behavior of logical circuits.
For a given circuit, the corresponding polynomial is obtained from the inputs and outputs of the circuit's gates as follows: 
\begin{itemize}
  \item A conjunction gate with inputs $a,b$ and output $s$ is represented by $-s + ab$
  \item A disjunction gate with inputs $a,b$ and output $s$ is represented by $-s + a + b - ab$
  \item An XOR-gate with inputs $a,b$ and output $s$ is represented by $-s + a + b - 2ab$
  \item A negation gate with input $a$ and output $s$ is represented by $-s + 1 -a$
\end{itemize}
Note that the roots of polynomials correspond to the semantics of the gates.
For a set of polynomials $P \subset \mathbb{Z}[X]$, that is obtained from the gates of a circuit, the collection of models of $P \cup B(X)$ is the set of all possible states of the circuit.
The set $P$ is called the \textit{gate constraints} of the given circuit.

Polynomials are also used to denote the \textit{specification} of a circuit.
For instance, for a multiplier circuit with inputs $a_0, \dots, a_{n-1}, b_0, \dots, b_{n-1}$ and outputs $s_0, \dots, s_{2n-1}$, the specification $f$ is given by
\begin{gather*}
  f := \sum_{0 \le i < 2n} 2^is_i - \left(\sum_{0 \le i < n} 2^ia_i\right) \cdot \left(\sum_{0 \le i < n} 2^ib_i\right).
\end{gather*}
A proof that a circuit with gate constraints $P$ adheres to a specification $f$ is obtained by showing that all models of $G := P \cup B(X)$ are also a model of $f$, which is done by proving that $f \in \langle G \rangle$.
Such a proof is obtained through stepwise derivation of polynomials in $\langle G \rangle$, such that the final step typically derives $f \in \langle G \rangle$.
Polynomials in $\langle G \rangle$ are obtained through linear combinations $p_1 \cdot q_1 + \dots + p_n \cdot q_n$, where for $1 \le i \le n$, $p_i \in \langle G \rangle$ and $q_i \in \mathbb{Z}[X]$.
Given a list of such derivations, a PAC checker checks the soundness of each derivation.
Note that, by construction, every model of $G$ is a root of every polynomial in $\langle G \rangle$.

Concretely, A PAC checker holds a state $(\mathscr{V},P)$, where $\mathscr{V}$ is a set of variables and $P$ is partial function $\mathbb{N} \rightharpoonup \mathbb{Z}[\mathscr{V}]$.
Initially it holds that $\mathscr{V} = X$ and $\text{img}(P)$ is the set of gate constraints, but throughout execution, further variables may be added to $\mathscr{V}$ and new polynomials (at a new index) are added to $P$.
The Boolean value constraints $B(X)$ are not in $\text{img}(P)$.
Instead, throughout the execution of the PAC checker, all monomials ${x_1}^{e_1} \dots {x_n}^{e_n} \cdot z$, for $x_1,\dots,x_n \in \mathscr{V}, e_1,\dots,e_n \in \mathbb{N}, z \in \mathbb{Z}$ are normalized to $x_1  \dots  x_n \cdot z$.
A PAC checker is given as input:
\begin{itemize}
  \item A list $[(i_1,p_1),\dots,(i_n,p_n)]$, where $i_1,\dots,i_n \in \mathbb{N}$ and $p_1, \dots, p_n \in \mathbb{Z}[X]$,
  \item a \textit{specification} $f \in \mathbb{Z}[X]$,
  \item and a list of \textit{proof steps}. A proof step is one of the following:
  \begin{itemize}
    \item A deletion step $\text{Del} (i)$, where $i \in \mathbb{N}$, or
    \item an extension step $\text{Ext} (i, x, p)$, where $i \in \mathbb{N}$, $x \not\in \mathscr{V}$ and $p \in \mathbb{Z}[\mathscr{V}]$, or
    \item a linear combination step $\text{LinComb} (i, r, [(i_1, q_1),  \dots, (i_n, q_n)])$, where
          $i \in \mathbb{N}$, $r \in  \mathbb{Z}[\mathscr{V}]$, and for all $1 \le j \le n$ $i_j \in \mathbb{N}$ and $q_j \in \mathbb{Z}[\mathscr{V}]$.
  \end{itemize}
\end{itemize}
The state of a PAC checker is initialized by $(\mathscr{V}_0 := \text{Var}(G \cup \{f\}), P_0 := \{i_j\mapsto p_j \mid 1 \le j \le n\})$.
For $i \not\in \text{dom}(P)$, I write $P(i) = \bot$.
I write $P(j\mapsto p)$ to denote $P$ with $j$ mapped to $p$.
Proof steps update the state $\sigma = (\mathscr{V},P)$ as follows:
\begin{itemize}
  \item For a deletion step $\text{Del} (i)$, the state is updated to $\sigma' := (\mathscr{V},P(i\mapsto \bot))$.
  \item For an extension step $\text{Ext} (i,x,p)$, the state is updated to $\sigma' := (\mathscr{V} \cup \{x\}, P(i \mapsto -x + p))$ if it holds that $P(i) = \bot$, $x \notin \mathscr{V}$ and $p \in \mathbb{Z}[X]\setminus B(X)$.
  \item For a linear combination step $\text{LinComb} (i,r,[(i_1,q_1),\dots,(i_n,q_n)])$, the state is updated to $\sigma' := (\mathscr{V}, P(i \mapsto r))$ if it holds that $P(i) = \bot$, $\forall j \in \{1,\dots,n\}. P(i_j) \neq \bot$ and $r = \sum_{1 \le j\le n}P(i_j)q_j \text{ mod } B(X)$.
\end{itemize}

\subsection{Isabelle/HOL}\label{sec:isabelle-hol}

The proof assistant \textit{Isabelle} \cite{isabelle} enables the writing of machine-checked formal mathematical proofs.
Following the LCF approach \cite{computational-logic-origins}, Isabelle employs a small kernel of logical inference rules for verifying the correctness of user-defined proofs, such that soundness of the kernel's inference engine implies correctness of the proofs.

To interactively construct proofs, the user can combine a variety of proof methods, ranging from fundamental deductive concepts, like case analysis and induction, to more complex automated methods.
An example of the latter is the \inlinecode{smt} method: A theorem is given to an external SMT solver, which emits a proof certificate.
The \inlinecode{smt} method then attempts to construct a proof from the certificate on the Isabelle side, using only existing Isabelle proof methods, yielding a valid proof that Isabelle's kernel accepts.
Crucially, errors in a proof method do not lead to acceptance of an unsound proof since the kernel rejects inference steps if they are not valid.

The core Isabelle system, referred to as \textit{Isabelle/Pure}, is a generic theorem proving framework, meaning it can serve as a basis for implementation of different logical theories with distinct axioms.
The most widespread instance of Isabelle is \textit{Isabelle/HOL}, which is an axiomatization of simply typed higher-order logic on top of Isabelle/Pure.
Throughout this thesis, I use Isabelle/HOL.

% Isabelle has been used for formalizations from various domains, including fundamentals of mathematics, e.g. Gödel's Incompleteness Theorems \cite{Incompleteness-AFP} or the Prime Number Theorem \cite{Prime-Number-Theorem-Avigad} \cite{Prime_Number_Theorem-AFP}, but also the verification of computer systems, for instance in the seL4 OS microkernel \cite{seL4}.

Throughout this thesis, I assume familiarity with simply typed higher-order logic and use notation closely tied to actual Isabelle syntax.
\begin{itemize}
  \item Leading apostrophes indicate type variables, e.g. $\tv{a}$.
  \item I use $\dbcolon$ to annotate the type of a value, e.g. $a \dbcolon \tv{a}$ means that value $a$ is of type $\tv{a}$.
  \item An algebraic datatype $\iexp{ta}$ with type parameters
        $\tv{a_1}, \dots, \tv{a_k}$ and $n$ constructors is denoted by
        $$
          (\tv{a_1}, \dots, \tv{a_k}) \app \iexp{ta}
            = C_1 \app \tau_{1,1} \app \dots \app \tau_{1,m_1}
            \mid \dots \mid
              C_n \app \tau_{n,1} \app \dots \app \tau_{n,m_n}
        $$
        where each constructor $C_i$ carries $m_i \ge 0$ arguments, and each
        argument type $\tau_{i,j}$ is an arbitrary type that may mention the
        type parameters $\tv{a_1}, \dots, \tv{a_k}$ as well as the defined type
        $(\tv{a_1}, \dots, \tv{a_k}) \app \iexp{ta}$ itself, allowing recursive
        datatypes.
  \item Functions are first-class objects and are typically presented in curried form.
        A function $f \dbcolon \tv{a} \rightarrow \tv{b} \rightarrow \tv{c}$, for a given $a \dbcolon \tv{a}$, returns a function of type $\tv{b} \rightarrow \tv{c}$.
  \item For a function $f \dbcolon \tv{a} \rightarrow \tv{b} \rightarrow \tv{c}$, $f\app a\app b$ denotes function application for arguments $a \dbcolon \tv{a}, b \dbcolon \tv{b}$.
  \item $\lambda$ denotes anonymous functions, e.g. $\lambda \app a \app b. f \app a \app b$ is an anonymous function that applies the named function $f$ to arguments $a$ and $b$.
  \item For types $\tv{a}$ and $\tv{b}$, $\tv{a} \times \tv{b}$ denotes the product type of tuples $(a \dbcolon \tv{a}, b \dbcolon \tv{b})$.
  \item For a type $\tv{a}$, $\tv{a} \app \iexp{list}$ denotes the type of lists with elements of type $\tv{a}$.
        A concrete list is denoted by $[x,y,z]$.\\
        The infix operator $\iexp{\#} \dbcolon \tv{a} \rightarrow \tv{a} \app \iexp{list} \rightarrow \tv{a} \app \iexp{list}$ prepends an element $x \dbcolon \tv{a}$ to a list $xs \dbcolon \tv{a} \app \iexp{list}$, denoted by $x \iexp{\#} xs$.\\
        The infix operator $\iexp{@} \dbcolon \tv{a} \app \iexp{list} \rightarrow \tv{a} \app \iexp{list} \rightarrow \tv{a} \app \iexp{list}$ concatenates two lists, denoted by $xs \iexp{@} ys$ for two lists $xs \dbcolon \tv{a} \app \iexp{list}$ and $ys \dbcolon \tv{a} \app \iexp{list}$.\\
        Isabelle/HOL defines the standard higher-order functions $\iexp{map} \dbcolon (\tv{a} \rightarrow \tv{b}) \rightarrow \tv{a} \app \iexp{list} \rightarrow \tv{b} \app \iexp{list}$ and $\iexp{foldl} \dbcolon (\tv{a} \rightarrow \tv{b} \rightarrow \tv{a}) \rightarrow \tv{a} \rightarrow \tv{b} \app \iexp{list} \rightarrow \tv{a}$.
        A right-fold also exists, but I do not use it in the thesis.
  \item For a type $\tv{a}$, $\tv{a} \app \iexp{set}$ denotes the type of sets with elements of type $\tv{a}$.
        $\{x,y,z\}$ denotes a concrete set.\\
        The function $\iexp{set} \dbcolon \tv{a} \app \iexp{list} \rightarrow \tv{a} \app \iexp{set}$ denotes the set of all elements contained in the given list.
  \item For a type $\tv{a}$, $\tv{a} \app \iexp{multiset}$ denotes the type of multisets with elements of type $\tv{a}$.
        $\{\#x,x,y\#\}$ denotes a concrete multiset.
        The function $\iexp{mset} \dbcolon \tv{a} \app \iexp{list} \rightarrow \tv{a} \app \iexp{multiset}$ denotes the multiset of all elements contained in the given list.
  \item For a type $\tv{a}$, $\tv{a} \app \iexp{option}$ denotes Isabelle's optional type
        $$
          \tv{a} \app \iexp{option} = \iexp{None} \mid \iexp{Some} \app \tv{a}
        $$
  \item I furthermore use the following built-in Isabelle types:
  \begin{itemize}
    \item Type $\iexp{bool}$ of booleans, with values $\iexp{True}$ and $\iexp{False}$,
    \item Type $\iexp{nat}$ of natural numbers (unbounded),
    \item Type $\iexp{int}$ of integers (unbounded),
    \item Type $w \app \iexp{word}$ of machine words with bit-width $w$,
    \item Unit type $\iexp{unit}$ with the only value $\iexp{()}$,
    \item Type $\iexp{char}$ of characters, where a character is defined by eight booleans, i.e. $\iexp{char} = \iexp{Char} \app (\iexp{digit0} \dbcolon \iexp{bool}) \app \dots \app (\iexp{digit7} \dbcolon \iexp{bool})$,
    \item Type $\iexp{string} = \iexp{char} \app \iexp{list}$.
  \end{itemize}
\end{itemize}

\subsection{Isabelle Refinement Framework}\label{sec:isabelle-refinement-framework}
The Isabelle Refinement Framework by Lammich \cite{isabelle-refinement-framework} enables a refinement approach to defining imperative programs and proving correctness properties of such programs. \todo{In motivation, write a section on what the idea of refinement is, and which problem it solves.}
This subsection introduces a simplified view of the underlying $\iexp{nres}$ monad and sketches its usage for defining and verifying imperative programs in Isabelle/HOL.
Most of the algorithms I implemented for this thesis are expressed in the $\iexp{nres}$ monad, the exception being fundamental operations on custom datatypes that \refsec{sec:low-level-datatypes-in-isabelle-llvm} introduces.

Nondeterministic computations are represented using the $\iexp{nres}$ monad.
$\iexp{nres}$ over some type $\tv{a}$, representing computations with results of type $\tv{a}$, is defined as\footnote{The presented $\iexp{nres}$ type is slightly simplified; the original version is defined in the theory $\iexp{Refine\_Monadic.Refine\_Basic}$.}
$$
  \tv{a} \app \iexp{nres} = \iexp{FAIL} \mid \iexp{RES} \app (\tv{a} \app \iexp{set})
$$
Intuitively, for $\iexp{RES} \app S$, the set $S$ represents all possible results of a computation, while $\iexp{FAIL}$ can be seen as a representation of failed computations, such as an out-of-bounds read from an array.

For lifting a value $a$ of type $\tv{a}$ into $\tv{a} \app \iexp{nres}$, the Refinement Framework defines
$$
  \iexp{RETURN} \dbcolon \tv{a} \rightarrow \tv{a} \app \iexp{nres}, \text{ where } \iexp{RETURN} \app a = \iexp{RES} \app \{a\}
$$
As an example, $\iexp{RETURN}$ is used to describe a program $\iexp{sum}$ that computes the sum of two natural numbers as follows:
$$
  \iexp{sum} \dbcolon \iexp{nat} \rightarrow \iexp{nat} \rightarrow \iexp{nat} \app \iexp{nres}, \text{ where } \iexp{sum} \app a \app b = \iexp{RETURN} \app (a \iexp{+} b)
$$
The Refinement Framework also defines $\iexp{SPEC}$, parameterized over a characteristic function $\Phi$ to describe results:
$$
  \iexp{SPEC} \dbcolon (\tv{a} \rightarrow \iexp{bool}) \rightarrow \tv{a} \app \iexp{nres}, \text{ where } \iexp{SPEC} \app \Phi = \iexp{RES} \app \{ a \dbcolon \tv{a} \mid \Phi \app a \}
$$
$\iexp{SPEC}$ is typically used to give specifications for programs.
For example, I define $\iexp{add\_spec}$ as a specification for a program that computes the sum of two natural numbers as follows:
$$
  \iexp{add\_spec} \dbcolon \iexp{nat} \rightarrow \iexp{nat} \rightarrow \iexp{nat} \app \iexp{nres}, \text{ where } \iexp{add\_spec} \app a \app b = \iexp{SPEC} \app (\lambda \app r.\, r = a \iexp{+} b)
$$
On $\iexp{nres}$, the Refinement Framework defines an order $\le$, given by
$$
\begin{aligned}
  \iexp{\_} &\le \iexp{FAIL} &&\Longleftrightarrow \iexp{True}, \text{ where $\iexp{\_}$ is an arbitrary value} \\
  \iexp{FAIL} &\le \iexp{RES} \app X &&\Longleftrightarrow \iexp{False} \\
  \iexp{RES} \app X &\le \iexp{RES} \app Y &&\Longleftrightarrow X \subseteq Y
\end{aligned}
$$
Intuitively, $P \le P'$ means that every possible result of program $P$ is also a possible result of $P'$.
A program $P \dbcolon \tv{a} \app \iexp{nres}$ is said to \textit{refine} program $P' \dbcolon \tv{a} \app \iexp{nres}$ if $P \le P'$.
In particular, if $P'$ is known to be correct, with respect to some definition of correctness, then $P \le P'$ implies that $P$ is also correct with respect to that definition of correctness.
For instance, $P \app a \app b \le \iexp{add\_spec} \app a \app b$ states that $P$ computes the sum of two natural numbers $a,b$.

Note that $(\tv{a} \app \iexp{nres}, \le)$ is a complete lattice with top element $\iexp{FAIL}$ and bottom element $\iexp{RES} \app \{\}$.
In particular, a refining program may only fail if the refined program also fails.

A $\iexp{bind} \dbcolon \tv{a} \app \iexp{nres} \rightarrow (\tv{a} \rightarrow \tv{b} \app \iexp{nres}) \rightarrow \tv{b} \app \iexp{nres}$ operator on $\iexp{nres}$, coupled with do-notation syntax, enables convenient denotation of sequential programs.
In addition, while loops and Isabelle's built-in $\iexp{if}$ expressions can be used to define programs in a typical imperative style.
\isabellecodeextern[linerange=59-74,label=algo:sum-imperative,caption={
  Imperative-style example program for computing the sum of a list of natural numbers.
}]{isabelle-code/LLVM_Basics.thy}
\refalgo{algo:sum-imperative} shows an example program for computing the sum of a list of natural numbers, using a while loop.
The function $\iexp{length} \dbcolon \tv{a} \app \iexp{list} \rightarrow \iexp{nat}$ returns the length of the given list.
The function $\iexp{!} \dbcolon \tv{a} \app \iexp{list} \rightarrow \iexp{nat} \rightarrow \tv{a}$ returns the element at a given position in the given list.
$\iexp{!}$ is used in infix notation, e.g. $xs \app \iexp{!} \app 3$ returns the fourth element of list $xs$, as indices start at $0$.
In the presented $\iexp{WHILET}$ notation, the first lambda function specifies the loop condition; if $i < \iexp{length} \app xs$ evaluates to $\iexp{False}$, the loop is exited.
The second lambda function then specifies the loop body, in which the accumulator and index are updated.

Correctness statements like $P \app a \app b \le \iexp{add\_spec} \app a \app b$ must be proven in Isabelle by the user.
The Refinement Framework offers tools to automate significant portions of such proofs.
In particular, the $\iexp{refine\_vcg}$ verification condition generator can be used to transform goals expressed in the $\iexp{nres}$ monad to standard Isabelle/HOL goals, on which Isabelle's proof automation (in particular $\iexp{sledgehammer}$) can then often be used effectively.

\subsection{Sepref}\label{sec:sepref}

In general, a program $\iexp{P} \dbcolon \tv{a} \app \iexp{nres}$ cannot be directly translated to an executable program, as it only describes the possible outcomes of a program.
This section gives an informal overview of the Sepref tool \cite{sepref} by Lammich, which is used to synthesize an executable implementation from such a program.

The Imperative-HOL framework \cite{Imperative-HOL} defines a semantics of imperative programs in Isabelle/HOL.
Programs defined in Imperative-HOL can then be exported to Standard-ML code.
Sepref utilizes a separation logic, using a heap monad in Imperative-HOL \cite{Separation_Logic_Imperative_HOL-AFP} which enables reasoning on imperative programs defined in Imperative-HOL.
Note that Imperative-HOL and the separation logic on it are not discussed here, since, for the thesis, I work with Isabelle-LLVM \cite{isabelle-llvm} (also by Lammich) which can be seen as Sepref's successor and defines its own separation logic on a shallow embedding of LLVM instead of Imperative-HOL (cf. \refsec{sec:isabelle-llvm}).

Consider the example program in \refalgo{algo:nres-set-lookup}, lifting a simple set membership test into $\iexp{nres}$.
\isabellecodeextern[linerange=31-32,label=algo:nres-set-lookup,caption={
  An example program in the $\iexp{nres}$ monad, checking whether a set contains a given element.
}]{isabelle-code/NRES_Fundamentals.thy}
As is, it is impossible to obtain an executable program for $\iexp{contains\_elem}$, since $\iexp{set}$ is the HOL type of (potentially infinite) sets, which cannot be represented with finite memory.
If the set $\iexp{s}$ is known to be finite, one could choose to implement $\iexp{contains\_elem}$ using an implementable datatype (e.g. a hashset) to represent $\iexp{s}$.
In such a case, the \textit{concrete} hashset datatype is said to \textit{refine} the \textit{abstract} set datatype.
\todo{Maybe a more rigorous introduction of datatype refinement here.}
The Sepref tool offers a collection of concrete data structures alongside refinement rules for implementing abstract datatypes.
In particular, Pastèque uses hashsets and hashmaps to refine sets and partial maps.

In addition, Sepref synthesizes Imperative-HOL programs and automates refinement proofs for given abstract programs, defined in the $\iexp{nres}$ monad. \todo{Should be more elaborate and mention the separation logic.}
For the synthesized Imperative-HOL programs, a code generator then generates Standard-ML code which, in the case of Pastèque, is compiled with MLton \cite{mlton,weeks-mlton-2006} to obtain an executable program.

In this scenario, the correctness of the synthesized program, relative to the semantics of Imperative-HOL and Sepref's separation logic on Imperative-HOL, is checked by the Isabelle kernel.
The compiler MLton and the code generator are trusted to be correct.

The newer Isabelle-LLVM framework uses the same $\iexp{nres}$ monad, but modifies the Sepref tool to target a shallow embedding of LLVM instead of Imperative-HOL, using its own separation logic.

Migrating the verified SAT-solver IsaSAT by Fleury \cite{isasat-sml} from the Standard-ML backend to LLVM \cite{isasat-llvm} improved its performance by a factor of two \todo{check number}.
This observation motivated the migration of Pastèque from Standard-ML to LLVM, which is the primary objective of this thesis.

\subsection{Isabelle-LLVM}\label{sec:isabelle-llvm}

This section gives an overview of the notion of \textit{refinement assertions} in Isabelle-LLVM and the separation logic used to reason about LLVM programs.

Isabelle-LLVM defines a shallow embedding of the LLVM intermediate language (i.e. a formalization of a subset of LLVM) and models the state of memory during the execution of a program.
Details regarding the LLVM embedding and memory model are described in Lammich's paper \cite{isabelle-llvm}.
The following definitions and terminology depend on the memory model and are only introduced informally here:
\begin{itemize}
  \item States of memory are reasoned about using a separation algebra $(\iexp{mem}, \#\#, +, 0)$ as defined in the AFP's Separation Algebra by Klein et al.~\cite{Separation_Algebra-AFP}, which is based on the Abstract Separation Logic by Calcagno et al.~\cite{Abstract_Separation_Logic}.
    Isabelle-LLVM instantiates the abstract separation algebra with its memory model.
  \item An \textit{assertion} is a function $f \dbcolon \iexp{mem} \rightarrow \iexp{bool}$ used to make statements about the state of memory.
    I use the type synonym $\iexp{assn} := \iexp{mem} \rightarrow \iexp{bool}$ as type annotation for assertions.
  \item $\iexp{emp}$ is an assertion that evaluates to $\iexp{True}$ iff. the given state is $0$, i.e. $\iexp{emp} = \lambda s. s = 0$.
  \item The separating conjunction operator $\ast \dbcolon \iexp{assn} \rightarrow \iexp{assn} \rightarrow \iexp{assn}$ allows for combination of assertions.
    If $A \dbcolon \iexp{assn}$ and $B \dbcolon \iexp{assn}$ are assertions, then $A \ast B$ is an assertion that evaluates to $\iexp{True}$ for a given state $h \dbcolon \iexp{mem}$, iff. there exists a decomposition $h = x + y$ such that $x \#\# y$ (i.e. $x$ and $y$ are disjoint) and $A \app x$ and $B \app y$.
  \item Some datatypes can be represented in LLVM.
        In particular, the type $w \app \iexp{word}$, of machine words with width $w$, where $w$ is one of $1,8,64,128$, is representable in LLVM and is used throughout the thesis.
        Note that other bit-widths are also representable, but are not used in the thesis.
  \item $\tv{a} \app \iexp{ptr}$ is the type of pointers to values of type $\tv{a}$.
        For $\mapsto \dbcolon \tv{a} \app \iexp{ptr} \rightarrow \tv{a} \rightarrow \iexp{assn}$, the assertion $(p \mapsto a) \sigma $ evaluates to $\iexp{True}$, iff. $p$ is a pointer to $a$ in state $\sigma$\footnote{$\mapsto$ is based on the assertion $\iexp{ll\_bpto}$, which is defined in the theory $\iexp{Isabelle\_LLVM.LLVM\_DS\_Block\_Alloc}$.}.
        Note that $p \mapsto a$ implicitly states exclusive ownership of the object $p$ points to, as no object can have multiple separate owners in the separation logic.
        \ref{sec:ownership-in-isabelle-llvm} discusses this more thoroughly as observation introduces the main challenge for the migration of Pastèque to Isabelle-LLVM.
  \item Similarly to the $\iexp{nres}$ monad, Isabelle-LLVM uses a state error nondeterminism monad $\tv{a} \app \iexp{llM}$ to describe sequential LLVM programs with return type $\tv{a}$.
        Like for $\iexp{nres}$, Isabelle-LLVM defines a bind operator with do-notation, loops and conditional statements to allow for expression of typical sequential imperative programs in $\iexp{llM}$. 
        Isabelle-LLVMs code generator can then generate actual LLVM code directly from programs defined using $\iexp{llM}$, given that all used types are representable in LLVM (alongside some other conditions which are not discussed here).
        While I don't give a formal introduction of the $\iexp{llM}$ monad, \refalgo{algo:sum-imperative-refined} shows an example program expressed in $\iexp{llM}$.
\end{itemize}

To make assertions independent of the state of memory, a boolean value can be lifted to an assertion by $\uparrow \dbcolon \iexp{bool} \rightarrow \iexp{assn}$ where $\uparrow \app b = \lambda h. \iexp{emp} h \land b$.

In Isabelle-LLVM, a \textit{refinement assertion} $A \dbcolon \tv{a} \rightarrow \tv{b} \rightarrow \iexp{assn}$ is a function that produces an assertion for given values of type $\tv{a}$ and $\tv{b}$, where type $\tv{b}$ is representable in LLVM.
Refinement assertions specify how abstract HOL types are represented in LLVM programs by concrete, implementable data types.
For example, a refinement assertion $\iexp{si8\_assn} \dbcolon \iexp{int} \rightarrow 8 \app \iexp{word} \rightarrow \iexp{assn}$ defines how 8-bit machine words, being interpreted as signed 8-bit integers, refine the abstract HOL type of unbounded integers.
It holds that $\iexp{si8\_assn} \app (-128) \app \iexp{0x80} \app 0 = \iexp{True}$, while $\iexp{si8\_assn} \app 128 \app w \app \sigma = \iexp{False}$ for any $w$ and $\sigma$, as $128$ is not representable with an 8-bit signed integer.
Throughout the thesis, for a refinement assertion $A \dbcolon \tv{a} \rightarrow \tv{b} \rightarrow \iexp{assn}$, I refer to $\tv{b}$ as the \textit{refinement target} of the abstract type $\tv{a}$.

To express correctness properties on programs defined in $\iexp{llM}$, Hoare-triples are used.
A \textit{Hoare-triple} $\{\Psi \dbcolon \iexp{assn}\} (P \dbcolon \tv{a} \app \iexp{llM}) \{\Phi \dbcolon \tv{a} \rightarrow \iexp{assn}\}$ states that, if a \textit{precondition} $\Psi$ holds initially, then the program $P$ is executed successfully and the \textit{postcondition} $\Phi$ holds afterwards.
For example, the Hoare-triple
$$
  \{\iexp{si8\_assn} \app a \app ai \ast \iexp{si8\_assn} \app b \app bi\}
  \app (\iexp{add} \app ai \app bi) \app
  \{\lambda r. \iexp{si8\_assn} \app (a + b) \app r\}
$$
can be read as follows:
If 8-bit machine words $ai, bi$ refine the integers $a, b$ then the result of $\iexp{add} \app ai \app bi$ refines the integer $a + b$.
Note that in this particular example, the statement only holds if $a + b$ does not overflow and is therefore representable by an 8-bit machine word.
This requirement can be added to the precondition:
\begin{gather*}
  \{\uparrow(a + b \le 127 \land -128 \le a + b) \ast \iexp{si8\_assn} \app a \app ai \ast \iexp{si8\_assn} \app b \app bi\} \\
  (\iexp{add} \app ai \app bi) \\
  \{\lambda \app r. \iexp{si8\_assn} \app (a + b) \app r\}
\end{gather*}
This addition makes the statement provable by the $\iexp{vcg}$ tool.

As a user of Isabelle-LLVM, one does not typically work with Hoare-triples directly.
Instead, the user can use the $\iexp{nres}$ monad to define algorithms using abstract HOL-types to then use the $\iexp{sepref}$ tool to transform the high-level programs into low-level programs expressed in the $\iexp{llM}$ monad and prove respective correctness statements automatically.
\isabellecodeextern[linerange=77-93,label=algo:sum-imperative-bound,caption={
  A variation of \refalgo{algo:sum-imperative}, where the length of the input list and the sum of the list are bound.
}]{isabelle-code/LLVM_Basics.thy}
Consider \refalgo{algo:sum-imperative-bound} which is the same program as in \refalgo{algo:sum-imperative} but with some additional conditions in the loop body, specified using $\iexp{ASSERT} \dbcolon \iexp{bool} \rightarrow \iexp{unit} \app \iexp{nres}$, where
\begin{gather*}
  \iexp{ASSERT} \app \varphi =
  \begin{cases}
    \iexp{RETURN} \app () &\text{, if $\varphi$} \\
    \iexp{FAIL} &\text{otherwise.}
  \end{cases}
\end{gather*}
Note that $\iexp{ASSERT}$ models assertions on $\iexp{nres}$ programs, analogously to $\iexp{assert}$ statements in programming languages like c or rust, and is not to be confused with the notion of assertions in separation logic, introduced in the beginning of this section.
The first assertion in the loop body of \refalgo{algo:sum-imperative-bound} ensures that the index $\iexp{i}$ is in bounds of the given list.
This assertion is mostly required for technical reasons, as it is directly implied by the loop guard.
The second assertion then ensures that the accumulator does not overflow, which also implies an upper bound on the sum of the elements in the given list.
With these assertions, the \inlinecode{sepref} tool can then synthesize a program $\iexp{sum\_impl} \dbcolon 64 \app \iexp{word} \times 64 \app \iexp{word} \app \iexp{ptr} \rightarrow 64 \app \iexp{word}$, shown in \refalgo{algo:sum-imperative-refined}.
\isabellecodeextern[linerange=101-112,label=algo:sum-imperative-refined,caption={
  $\iexp{llM}$-level implementation of \refalgo{algo:sum-imperative-bound}, synthesized by \inlinecode{sepref}.
}]{isabelle-code/LLVM_Basics.thy}
The type signature of $\iexp{sum\_impl}$ implies the concrete datatypes used: The list is implemented by an array of 64-bit (signed) integers, the sum is implemented by a 64-bit (signed) integer.
In the tuple $64 \app \iexp{word} \times 64 \app \iexp{word} \app \iexp{ptr}$, the first element saves the length of the array, while the second element is a pointer to the actual data. 
The concrete datatype is a choice the user must make.
In this example, one could also choose to, e.g., implement the abstract list by a linked list of 8-bit unsigned integers.

\inlinecode{sepref} translates the abstract function $\iexp{length} \dbcolon \tv{a} \app \iexp{list} \rightarrow \iexp{nat}$ and $\iexp{!} \dbcolon \tv{a} \app \iexp{list} \rightarrow \iexp{nat} \rightarrow \tv{a}$ to concrete functions $\iexp{la\_length\_impl}$ and $\iexp{la\_get\_impl}$ which return then length of, and a specific element from an array.
For both these functions, the array implementation also specifies corresponding correctness lemmas which are required for proving the overall correctness of $\iexp{sum\_impl}$.

Note that the original $\iexp{sum}$ program from \refalgo{algo:sum-imperative} could not be synthesized by \inlinecode{sepref} as in the concrete implementation $\iexp{sum\_impl}$, the accumulator could overflow, leading to a different end result compared to the abstract $\iexp{sum}$ which uses unbounded natural numbers.
More concretely, in addition to synthesized $\iexp{llM}$ programs, \inlinecode{sepref} also automatically proves refinement theorems.
For the $\iexp{sum\_impl}$ example, a simplified lemma using Hoare-triple notation of the corresponding refinement is given by
\begin{gather*}
  \{(\iexp{array\_assn} \app \iexp{sn64\_assn}) \app xs \app xsi\}\\
  (\iexp{sum\_impl} \app xsi) \\
  \{\lambda r. \app \iexp{sn64\_assn} \app (\iexp{sum'} \app xs) \app r\}
\end{gather*}
Where $\iexp{array\_assn} \app \iexp{sn64\_assn} \dbcolon \iexp{nat} \app \iexp{list} \rightarrow 64 \app \iexp{word} \times 64 \app \iexp{word} \app \iexp{ptr} \rightarrow \iexp{assn}$ is a refinement assertions specifying refinement of lists of natural numbers by arrays of 64-bit signed integers and $\iexp{sn64\_assn} \dbcolon \iexp{nat} \rightarrow 64 \app \iexp{word} \app \iexp{assn}$ specifies refinement of natural numbers by signed 64-bit integers.

As a final step, the $\iexp{llM}$-level program $\iexp{sum\_impl}$ in \refalgo{algo:sum-imperative-refined} can then be exported by Isabelle-LLVMs code generator to obtain compilable LLVM code.
The code generator can also generates header files so the exported functions can be called from c programs.
For Pastèque in particular, I have implemented the initial parsing of input files in c to then use the parsed inputs as arguments for the verified LLVM algorithms. 

\subsubsection{Ownership in Isabelle-LLVM}\label{sec:ownership-in-isabelle-llvm}
This subsection introduces the notion of \textit{pure} and \textit{impure} refinement assertions and how Isabelle-LLVMs backend is different to the Imperative-HOL backend with regards to the purity of lists.

Refinement Assertions are distinguished between \textit{pure} and \textit{impure}.
A refinement assertion $A \dbcolon \iexp{assn}$ is pure if it holds that $\forall s. \app A \app s \Rightarrow s = 0$.
Refinement assertions with stack allocated refinement targets, in particular machine words, are pure.
The introduced refinement assertions $\iexp{si8\_assn}$ and $\iexp{sn64\_assn}$ are therefore examples of pure refinement assertions.
Refinement assertions with heap allocated refinement targets, in particular data structures like arrays or lists, are impure.
$\iexp{array\_assn} \app \iexp{sn64\_assn}$ is an example of an impure refinement assertion.

Stack allocated values can be duplicated without requiring a specific copy implementation.
For example, the earlier Hoare-triple from \refsec{sec:isabelle-llvm} can be strengthened to
\begin{gather*}
  \{\uparrow(a + b \le 127 \land -128 \le a + b) \ast \iexp{si8\_assn} \app a \app ai \ast \iexp{si8\_assn} \app b \app bi\} \\
  (\iexp{add} \app ai \app bi) \\
  \{\lambda \app r.\, \iexp{si8\_assn} \app (a + b) \app r \ast \iexp{si8\_assn} \app a \app ai \ast \iexp{si8\_assn} \app b \app bi\}
\end{gather*}
The addition of $\iexp{si8\_assn} \app a \app ai \ast \iexp{si8\_assn} \app b \app bi$ means that, after being used as arguments for $\iexp{add}$, $ai$ and $bi$ remain valid representations of the original abstract integers $a$ and $b$ and could be used again in further computations.

Impure assertions behave differently with regards to preservation of values.
As an example, consider the function $\iexp{take} \dbcolon \iexp{nat} \rightarrow \tv{a} \app \iexp{list} \rightarrow \tv{a} \app \iexp{list}$ which returns a prefix of the given length for the given list.
Assuming a list of natural numbers $\iexp{nat} \app \iexp{list}$ is refined by an array of (signed) 64-bit machine words, i.e. $64 \app \iexp{word} \times 64 \app \iexp{word} \app \iexp{ptr}$, like in the $\iexp{sum\_impl}$ example (cf. \refalgo{algo:sum-imperative-refined}), one might want to prove a Hoare-triple
\begin{gather*}
  \{(\iexp{array\_assn} \app \iexp{si64\_assn}) \app xs \app xsi \app \ast \app \iexp{si64\_assn} \app i \app ii\} \\
  (\iexp{take\_impl} \app ii \app xsi) \\
  \{\lambda r. (\iexp{array\_assn} \app \iexp{si64\_assn}) \app (\iexp{take} \app i \app xs) \app r \app \ast \app (\iexp{array\_assn} \app \iexp{si64\_assn}) \app xs \app xsi\}
\end{gather*}
for an implementation $\iexp{take\_impl}$ of $\iexp{take}$ such that the original array $xsi$ is kept intact after executing $\iexp{take\_impl} \app ii \app xsi$.

In a language like c, one might implement $\iexp{take\_impl}$ in a way such that $xsi$ and $r$ both point to the same data in memory, but $r$ has length $ii$, while $xsi$ still has it's original length.
While this is efficient and has constant runtime complexity, in Isabelle-LLVM its impossible to implement $\iexp{take\_impl}$ this way since two distinct pointers must not point to the same object in memory.
In other words, each resource must have one unique owner.
\footnote{In recent developments, Isabelle-LLVM was extended to support fractional separation logic \cite{isabelle-llvm-fractional} which enables shared read-only access to an object and in this example might allow for using a shared, read-only prefix of $xsi$, but at the time of starting the work on this thesis, this work was not yet published.}
The actual implementation of the $\iexp{take}$ operation in Isabelle-LLVM therefore proves a different Hoare-triple
\begin{gather*}
  \{(\iexp{array\_assn} \app \iexp{si64\_assn}) \app xs \app xsi \app \ast \app \iexp{si64\_assn} \app i \app ii\} \\
  (\iexp{take\_impl} \app ii \app xsi) \\
  \{\lambda r. (\iexp{array\_assn} \app \iexp{si64\_assn}) \app (\iexp{take} \app i \app xs) \app r\}
\end{gather*}
Where in the postcondition the second conjunct is lost, meaning that the original array $xsi$ is destroyed and can no longer be used.
If $xsi$ must be kept-intact, $\iexp{take\_impl}$ would have to create a deep copy of the respective prefix, leading to linear runtime and allocation of a new array.

The Imperative-HOL backend, the refinement assertion $\iexp{list\_assn}$ on lists is actually a pure assertion.
Efficient handling of lists and creating deep copies where necessary is handled by the compiler and garbage collector.
In LLVM, this is not the case; lists are an impure chain of lists, meaning that any algorithm which uses functions like $\iexp{take}$ necessarily destroys given input lists.

As lists are frequently used in the implementation of Pastèque, this complicates the migration to Isabelle-LLVM significantly, unless every function working on lists creates a deep copy first.
\refsec{sec:lists-with-tail-pointers} introduces some measures I took to avoid deep copies by reformulating functions using $\iexp{foldl}$ where possible.

\subsection{Pastèque}\label{sec:methodology_high-level-representation-of-polynomials}
This subsection introduces the fundamental algorithms and HOL-level data types used in Pastèque, implemented by Fleury \cite{pasteque} \cite{pac-2024}.
The discussed implementations are done on $\iexp{nres}$ level and remain unmodified in the migration to LLVM, meaning that all my contributions are refinements of the existing algorithms (apart from the unverified parser which is an entirely new implementation in c).

To represent polynomials, Pastèque defines
\begin{align*}
  \iexp{variable}   &= \iexp{string} \\
  \iexp{term}       &= \iexp{variable} \app \iexp{list} \\
  \iexp{monomial}   &= \iexp{term} \times \iexp{int} \\
  \iexp{polynomial} &= \iexp{monomial} \app \iexp{list}
\end{align*}
For instance, a polynomial $3xy + 2x + y$ is represented by
$[([\iexp{"x"}], 2), ([\iexp{"x"}, \iexp{"y"}], 3), ([\iexp{"y"}], 1)]$ in Pastèque.

The map of polynomials $P$, as defined in \ref{sec:practical-algebraic-calculus}, is represented by a partial map $\iexp{A} \dbcolon \iexp{nat} \rightharpoonup \iexp{polynomial}$.\footnote{Partial maps $m \dbcolon \tv{a} \rightharpoonup \tv{b}$ in Isabelle/HOL are total functions $f_m \dbcolon \tv{a} \rightarrow \tv{b} \app \iexp{option}$, where $\forall a \dbcolon \tv{a}. \app a \in \iexp{dom} \app m \Rightarrow \exists b \dbcolon \tv{b}. f_m \app a = \iexp{Some} \app b$, while $\forall a \dbcolon \tv{a}. \app a \not\in \iexp{dom} \app m. f_m \app a = \iexp{None}$.}
The specification, called $f$ in \ref{sec:practical-algebraic-calculus}, is represented by a polynomial $\iexp{spec} \dbcolon \iexp{polynomial}$.
Proof steps are represented through an algebraic datatype
\begin{align*}
  \iexp{pac\_step} =  &\iexp{LinComb} \app ((\iexp{polynomial} \times \iexp{nat}) \app \iexp{list}) \app \iexp{nat} \app \iexp{polynomial} \\
                      &\mid \iexp{Ext} \app \iexp{nat} \app \iexp{string} \app \iexp{polynomial} \\
                      &\mid \iexp{Del} \app \iexp{nat}
\end{align*}
which directly models the notion of proof steps from \ref{sec:practical-algebraic-calculus}.

Throughout execution, all polynomials are sorted with respect to a lexicographic order on the monomial's terms, allowing for an efficient addition of polynomials.
As an example, consider two unsorted polynomials
\begin{align*}
  p_1 &:= 3x + 2y + z\\
  p_2 &:= 3z + 2y + x
\end{align*}
a naive implementation of addition of polynomials might concatenate $p_1$ and $p_2$, sort the result (with respect to the monomial's terms) and then merge monmials $3x,x \rightsquigarrow 4x$, $2y,2y \rightsquigarrow 4y$ and $z,3z \rightsquigarrow 4z$, amounting to an implementation for addition with runtime complexity $\mathcal{O}(n\log(n))$.
If polynomials are sorted however, addition is possible in linear time.
Consider two polynomials $p = c_{p1}m_{p1} + \dots c_{pn}m_{pn}, q = c_{q1}m_{q1} + \dots + c_{qk}m_{qk}$ that are both sorted with respect to their monomial's terms, i.e., $m_{p1} < \dots < m_{pn}$ and $m_{q1} < \dots < m_{qk}$.
The addition $r := p+q$ is computed by initializing $r=0$ and iterating over pairs $(c_{pj}m_{pj},c_{ql}m_{ql})$ of monomials from $p$ and $q$ as follows:
\begin{itemize}
  \item If $c_{pj} < c_{ql}$ then update $r' := r + c_{pj}m_{pj}$ and continue with the next pair of monomials $(c_{p(j+1)}m_{p(j+1)},q_{ql}m_{ql})$.
  \item If $c_{pj} > c_{ql}$ then update $r' := r + c_{ql}m_{ql}$ and continue with $(c_{pj}m_{pj},q_{q(l+1)}m_{q(l+1)})$.
  \item If $c_{pj} = c_{ql}$ then compute $s := m_{pj} + m_{ql}$.
        If $s \neq 0$, then update $r' := r + c_{pj}s$.
        Otherwise, do not update $r$.
        Proceed with $(c_{p(j+1)}m_{p(j+1)},c_{q(l+1)}m_{q(l+1)})$.
\end{itemize}
This is done until the end of either $p$ or $q$ is reached.
The remaining monomials of the other polynomial are appended to the result.
Note that the result $r$ is also sorted with respect to its monomials terms.

The current state in Pastèque is represented by a triple $(\mathscr{V}, \iexp{A}, \iexp{s}) \dbcolon \iexp{variable} \app \iexp{set} \times (\iexp{nat} \rightharpoonup \iexp{polynomial}) \times \iexp{status}$, where $\iexp{status} = \iexp{FAILED} \mid \iexp{SUCCESS} \mid \iexp{FOUND}$, where $\iexp{SUCCESS}$, which is also the initial state, is a kept as long as proof steps are valid, $\iexp{FAILED}$ is entered once an invalid proof step occurs and $\iexp{FOUND}$ is entered if all steps are valid and the target polynomial $\iexp{spec}$ is the result of a linear combination proof step.
Once all proofs steps were checked, or if $\iexp{FAILED}$ is entered, Pastèque returns its last state.

As an optimization to the $\iexp{monomial}$ type above, a variant of Pastèque uses \textit{shared variables} given by two mappings $\iexp{NV} \dbcolon \iexp{nat} \rightharpoonup \iexp{variable}$ and $\iexp{VN} \dbcolon \iexp{variable} \rightharpoonup \iexp{nat}$ such that $\iexp{NV} \circ {VN} = \iexp{id}$ and $\iexp{VN} \circ \iexp{NV} = \iexp{id}$, where $\circ$ is the composition of two functions and $\iexp{id}$ is the identity function.
This allows for a more efficient representation of polynomials given by $\iexp{sllist\_polynomial} :\equiv (\iexp{nat} \app \iexp{list} \times \iexp{int}) \app \iexp{list}$.
For the remainder of the thesis, I will refer to the $\iexp{string}$ based polynomials as $\iexp{llist\_polynomial}  :\equiv \iexp{polynomial}$ which is the nomenclature in the actual Pastèque source code.
With the shared variable maps, the state of Pastèque extends to $(\mathscr{V}, \iexp{A}, \iexp{s}, \iexp{NV}, \iexp{VN})$.
For the migration to Isabelle-LLVM, a refinement target for this state, which is representable in LLVM, is required.
Leading to the two main tasks for this thesis:
\begin{enumerate}
  \item The Isabelle/HOL $\iexp{int}$ type is unbounded and can hold arbitrarily large numbers.
  The Imperative-HOL backend, which produces SML code that is compiled by MLton, supports this natively.
  LLVM only supports bounded numbers, represented as machine words of a fixed width.
  Eberl et al. have used bindings to GMP \cite{gmp} in Isabelle-LLVM to implement a verified algorithm for computing Bernoulli numbers \cite{eberl-bernoulli-numbers}.
  Instead of GMP, I use a native library for arbitrary precision numbers, created by Spinei \todo{How can we cite Mihai?}.
  This library implements addition, subtraction and multiplication for unsigned numbers and can thus refine these operations on natural numbers.
  To use it for Pastèque, I extend the implementation to support signed numbers to refine integers used in Pastèque.
  Furthermore, I implement conversion between numbers and strings which is required for printing error messages and parsing input files.
  \refsec{sec:multiple-precision-integers} discusses the work on arbitrary-precision integers.
  \item The state representation defined above implies a need for container data types which are not directly available in Isabelle-LLVM:
  \begin{itemize}
    \item I extend the existing linked list implementation to obtain a nestable list implementation which is discussed in \ref{sec:lists-with-tail-pointers}.
    \item I compose and extend existing data structures to implement a partial map, yielding a refinement target for the map of stored polynomials $\iexp{A}$.
          The partial map implementation is discussed in \ref{sec:partial-maps}
    \item To refine the set of variables $\mathscr{V}$, I implement hash sets and define a simple hash function for strings.
    \item I refine $\iexp{VN}$ by a custom hash map implementation, reusing the hash function from the hash set of strings.
  \end{itemize}
\end{enumerate}
